\section{Conclusion}
In this paper, we propose FreeArtGS, a novel method for reconstructing free-moving articulated objects from monocular RGB-D videos. Our method first segments the free-moving parts by combining an optimization-based method with point-tracking priors. Based on the estimated part segments and transformations, it then infers the joint type and axis by fitting the relative motion between parts. Finally, a 3DGS-based end-to-end optimization jointly refines the joint parameters, geometry, and appearance. Experiments demonstrate the robustness and effectiveness of our method in both simulated and real-world settings. With the growing need to rapidly expand articulated digital twins for augmented reality and robotics, our method provides a promising solution with fewer constraints and higher scalability.

FreeArtGS still has several limitations. First, our method currently assumes a two-part articulated object; extending it to multi-part structures by sequentially capturing each moving part remains an important direction. Second, relying on multiple off-the-shelf priors can lead to cascading error accumulation. A potential solution lies in developing a unified feed-forward model to simultaneously predict joints, poses, geometry, and textures. Third, our framework requires monocular RGB-D input. While extending it to RGB-only sequences by predicting continuous video depth is a natural progression, it currently faces challenges regarding depth accuracy. We leave these as future work.

% First, our method uses a two-part assumption, which needs further extension when applied to multi-part articulated objects. 
% Second, it leverages the priors from off-the-shelf models as initializations or supervisions. Although we have designed
% optimization-based noise resistance techniques to prevent
% the prediction errors, the whole pipeline still needs careful parameter tuning. 
% Third, due to the setting of monocular RGB-D video, in further real-world applications, our method can be influenced by a larger range of depth noise.
